% Source: https://artofproblemsolving.com/community/c3h3299913_regarding_irrational_numbers
% Stable identity: aops-topic-3299913-q01
(a) If two decimal expansions first differ at position $k$, the largest possible contribution of all later differences is $9\sum_{j>k}10^{-j}=10^{-k}$. Equality requires one tail to be all $9$s and the other all zeroes. Thus two expansions of the same number can differ only in the terminating-versus-repeating-$9$ case, which is rational. An irrational number has no such ambiguity.
(b) Equal fractional parts would give $\sqrt m-\sqrt n=h$ for an integer $h\ne0$. Squaring yields $m=n+2h\sqrt n+h^2$, hence $\sqrt n=\dfrac{m-n-h^2}{2h}$ rational, contradicting that $n$ is nonsquare.
